\section{Introduction}
\label{sec:introduction}

Forecasting representations are commonly assessed through the errors of fitted
decoders. Such comparisons do not determine whether another decoder could make
a simpler representation meet the same accuracy requirement. We study how to
certify that distinction. An upper risk bound establishes the accuracy of a
fitted forecast; a lower risk bound can exclude every decoder in a declared
representation class. Together they can establish when a more costly
representation is necessary. All comparisons use the same observed target
and evaluation law, with decoder sensitivity budgets declared explicitly.
The added conclusion is necessity within a declared family: the least costly
accepted fit is certified to have minimum feasible cost when every cheaper
class is excluded. This conclusion can hold even when upper-only selection returns
the same fitted forecast.

Kernel Affine Hull Machines (KAHMs) provide our main worked example.
Their established local reconstruction models map histories to soft
associations with reference groups
\citep{kumar2024kahm,kumar2025collaborative,kumar2026limits}.
We retain these models and append feature blocks under a sum-of-block-norms
decoder budget. This preserves earlier feasible forecasts and makes exact
block duplication immaterial to the class. The problem is to assess whether
these learned coordinates support a forecasting tolerance, and whether an
added block is necessary to attain it. The certificates apply to other fixed
representations as well; no theorem assumes that reconstruction affinity
establishes predictive sufficiency or exact linear feature evolution.

A basic obstruction arises when two targets are farther apart than the
sensitivity budget permits their predictions to be. Related pairwise
Lipschitz regression arguments are established \citep{akash2021inverse}.
Averaging the squared positive part of this discrepancy over conditionally
independent pairs yields a risk lower bound. This inexpensive statistic can
be exact, but it can also miss almost all of the attainable class risk.
We ask how to reduce that population slack while accounting separately for
uncertainty from a finite calibration sample.

Our main construction replaces pairs by blocks of observations. Within each
block, we minimize squared error over predictions satisfying all pairwise
Lipschitz constraints and normalize by one less than the block size. The
size-two statistic is exactly the pairwise obstruction. For scalar targets,
the expectations of these statistics form a nondecreasing hierarchy of lower
bounds. Under unconditional iid sampling, bounded targets, and compact
representation support, they converge to the optimal Lipschitz risk. Larger blocks capture compatibility constraints that
individual pairs cannot express. They also require more optimization and
leave fewer independent block scores for calibration. Population improvement
therefore does not imply a uniformly stronger finite-sample certificate.
Monotonic lower expectations of ordinary sample-average minima are established
\citep{mak1999montecarlo}. Here the centered normalization connects the
Lipschitz problem exactly to the pair score, and the rare-conflict and
noisy-support laws quantify the information missed by finite blocks.

Conditional pairing offers a complementary approach. Comparisons within a
predefined history stratum can expose conflicts that global pairing rarely
encounters. Reweighting by stratum probabilities preserves the marginal risk;
calibration accounts for random counts and estimated masses. A lower bound
nevertheless assesses the entire representation--decoder class: positive risk
may reflect information loss, irreducible future noise, or the sensitivity
constraint. We do not identify these causes from a lower endpoint alone.

The paper makes three contributions.
\begin{enumerate}[leftmargin=*,itemsep=2pt]
\item \textbf{A centered hierarchy and its population limits.} We connect the pair
witness to normalized minimum-risk problems on larger blocks, establish
monotonicity and consistency under the stated conditions, and identify exact
and arbitrarily loose pairwise cases.
\item \textbf{Representation decisions from lower and upper bounds.}
Independent block scores and stratification give class-wide lower
certificates. Combined with fitted-risk upper bounds and retained-block
geometry, they support accuracy, improvement, and minimum feasible cost
statements for a frozen family. The concentration inequalities and
simultaneous-selection principles are established components.
\item \textbf{Tests of usefulness and failure.} Exact-risk comparisons separate
population slack, precision, and optimization cost. Fresh calibration on 40
frozen KAHM feature banks shows larger blocks changing class-exclusion
decisions, while retaining cases where full-sample empirical minimization is
stronger. Controlled refinement and broader synthetic, pen-trajectory, and
speech studies assess accuracy, scaling, and calibration limits.
\end{enumerate}

The evidence separates informative decisions from unresolved ones.
Empirical-minimum calibration is often tighter than pairwise calibration,
and KAHM accuracy depends on the task and decoder. Useful class exclusion
on a constructed mechanism does not establish its prevalence on natural data.
All confidence statements concern expected squared error under the declared
sampling law. Sections~\ref{sec:theory}--\ref{sec:block-witness} develop the
population bounds; Sections~\ref{sec:method}--\ref{sec:calibration} connect them
to fixed forecasting procedures and calibrated decisions.
